% ============================================================
% 付録：AI活用記録の書き方、プロンプトの4要素、出力を使う前のチェック
% ============================================================
\colorlet{phase}{mone}\colorlet{phasetint}{monetint}
\intro{付録：AI活用記録とプロンプトの基本}

\topic{AI活用記録（提出物に添付）}
提出物には、次の形式で「AI活用記録」を添えてください。長くなくてかまいません。使っていない場合は「使用なし」と書きます。
\par\noindent\begin{gtab}{colspec={Q[l,wd=34mm] X[l]}}
  使用したツール & Copilot（大学提供） \\
  使用した場面 & 例：問いの候補出し、検索キーワードの検討、要約、反論の生成、表現の推敲 \\
  主なプロンプト & 実際に使った指示を1〜3個。長い場合は要点でよい \\
  出力をどう扱ったか & 採用した点、修正した点、採用しなかった点とその理由 \\
  検証したこと & 事実、文献、数値を、何で確認したか \\
\end{gtab}
レポート本文の冒頭または末尾には、次のような一文を入れます。
\begin{promptbox}[本文での明記]
本レポートの作成にあたり、問いの候補出しと構成案の検討に大学が提供する Copilot を使用した。出力はすべて筆者が一次資料に基づいて検証し、文章は筆者自身が執筆した。
\end{promptbox}
AIの出力そのものを引用する必要がある場合は、「Copilot に『〜』と入力した際の出力（2026年10月7日）」のように、プロンプトと日付がわかる形で示します。

\topic{プロンプトの4要素}
AIとの対話の質は、指示の質で決まります。本講義では次の4つの要素を意識します。
\par\noindent\begin{gtab}{colspec={Q[l,wd=28mm] X[l]}}
  役割 & AIに立ってほしい立場。「あなたは経営学を専門とする大学教員です」「あなたは厳しい査読者です」 \\
  目的と背景 & 何のために、どんな状況で必要か。「経営情報の授業で、企業へのAI活用提案を作っています」 \\
  条件 & 範囲、制約、避けてほしいこと。「日本国内の事例に限る」「文献名は挙げないでください」 \\
  出力の形式 & 箇条書き、表、文字数、案の数。「300字以内」「10案を表で」 \\
\end{gtab}
対話を深めるコツは、一度で終わらせないことです。出力に対して「なぜ」「他には」「反対の立場だと」と続け、
良い出力があっても「この案の弱点は」と聞いてから採用します。自分の考えを先に書いてからAIに批評させると、
AIの案に引きずられにくくなります。

\topic{出力を使う前のチェック}
\begin{checks}
  \item 固有名詞（人名、組織名、書名、論文名）は実在するか。
  \item 数値には出典があるか。出典にあたって同じ数値か。
  \item 「〜と言われている」「一般に」など、根拠のない断定になっていないか。
  \item 反対の見解や例外はないか。「この主張に反する証拠は」と聞いたか。
  \item 自分で説明できるか。説明できない内容は使わない。
\end{checks}

\topic{関連する教材}
\begin{itemize}
  \item 講義サイトの「AI利用ガイド」：本講義でのAIの使い方の原典。してよいこと・気をつけること・しないことの一覧。
  \item AI利用ガイドの復習教材（Teams で配布）：ハルシネーションの確認、問いの立て方、先行研究の探し方、文献の批判的な読み方、問いを磨く、の5単元。
  \item 講義サイトの「プログラミング入門」：HTML・CSS・JavaScript・Python を1日1テーマで学ぶ30日の自習教材。第8〜10回でデータを扱う際の基礎になる。
\end{itemize}
